% =========================================================================
% Carga y configuración de paquetes, agrupados por función y con el porqué
% de cada uno. hyperref se carga casi al final (requisito del paquete).
% =========================================================================

% --- Codificación y fuentes (pdflatex) ---
\usepackage[utf8]{inputenc}   % Codificación de entrada UTF-8
\usepackage[T1]{fontenc}      % Acentos correctos en el PDF (copiables/buscables)

% --- Idiomas: español principal; euskera e inglés para el resumen trilingüe ---
% es-tabla: "Tabla" en vez de "Cuadro".
% es-noshorthands: desactiva los caracteres activos del español (no los
%   necesitamos: escribimos los acentos en UTF-8) y evita un conflicto real
%   entre babel-spanish y el hook de biblatex-apa sobre \@makecaption
%   ("Generic hooks cannot be added to '\@makecaption'").
\usepackage[english, basque, spanish, es-tabla, es-noshorthands]{babel}
\usepackage{csquotes}         % Comillas correctas según idioma (lo pide biblatex)
% Nomenclatura de la guía de la EIB: "Anexo" para los apéndices e "Índice de
% ilustraciones" para las figuras ("Índice de tablas" ya es el defecto con
% es-tabla):
\addto\captionsspanish{%
  \renewcommand{\appendixname}{Anexo}%
  \renewcommand{\listfigurename}{Índice de ilustraciones}%
}

% --- Gráficos y tablas ---
\usepackage{graphicx}
\graphicspath{{figures/}}     % Las imágenes se buscan siempre en figures/
\usepackage{booktabs}         % Tablas profesionales (\toprule, \midrule, \bottomrule)

% --- Utilidades ---
\usepackage{xspace}           % Espaciado correcto tras macros de texto (\bukit)
\usepackage{xcolor}           % Colores (usado por \todonote)
\usepackage{acronym}          % Lista de acrónimos (ver frontmatter/acronyms.tex)

% --- Bibliografía: biblatex + biber, en normas APA (exigido por la guía) ---
% Citar con \parencite{clave} -> "(Williams, 1995)"
%       o  \textcite{clave}   -> "Williams (1995)"
\usepackage[backend=biber, style=apa]{biblatex}
\DeclareLanguageMapping{spanish}{spanish-apa}
% biblatex-apa no trae localización en euskera (basque-apa.lbx); mapeamos a
% la española para silenciar el aviso. Solo afectaría a citas hechas DENTRO
% del bloque en euskera (el Laburpena), donde nunca citamos.
\DeclareLanguageMapping{basque}{spanish-apa}

% --- Enlaces y metadatos del PDF (cargar al final) ---
\usepackage[hidelinks]{hyperref}
\hypersetup{
  pdftitle={\thesisTitle},
  pdfauthor={\thesisAuthor},
  pdfsubject={\thesisSubtitle},
  pdfkeywords={\thesisKeywordsEs},
}
